\documentclass[11pt,a4paper]{article}
\PassOptionsToPackage{pdfusetitle=false}{hyperref}
\usepackage[utf8]{inputenc}
\usepackage[T1]{fontenc}
\usepackage[margin=2.5cm]{geometry}
\usepackage{amsmath,amssymb}
\usepackage{booktabs,array}
\usepackage{xcolor}
\usepackage{tcolorbox}
\tcbuselibrary{skins,breakable}
\usepackage{enumitem}
\usepackage[pdfusetitle=false]{hyperref}
\usepackage[numbers]{natbib}
\bibliographystyle{unsrtnat}
\hypersetup{colorlinks=true,linkcolor=blue!60!black,urlcolor=blue!60!black,
            citecolor=blue!60!black,
            pdftitle={Particles as Topological Kinks},
            pdfauthor={Nova Spivack}}
\setlength{\emergencystretch}{2em}
\usepackage{tikz}
\usetikzlibrary{arrows.meta,positioning,shapes.geometric,calc,fit,backgrounds}

%─── Color palette ───────────────────────────────────────────────────────────
\definecolor{boxblue}{RGB}{220,235,255}
\definecolor{boxorange}{RGB}{255,240,210}
\definecolor{boxgreen}{RGB}{220,245,220}
\definecolor{boxtan}{RGB}{255,245,220}
\definecolor{boxpurple}{RGB}{240,225,255}
\definecolor{boxred}{RGB}{255,230,225}

%─── Macros ──────────────────────────────────────────────────────────────────
\newcommand{\term}[1]{\textbf{#1}}
\newcommand{\lean}[1]{%
  \penalty-100
  \begingroup\def\_{\textunderscore\allowbreak}\texttt{#1}\endgroup}
\newcommand{\CatAL}{\textbf{CatAL}}
\newcommand{\CatAD}{\textbf{CatAD}}
\newcommand{\CatA}{\textbf{CatA}}
\newcommand{\PHIMDL}{$\Phi_{\rm MDL}$}
\newcommand{\PhiMDL}{\Phi_{\rm MDL}}
\newcommand{\Mkink}{M_{\rm kink}}

\newcommand{\TutorialDOI}{10.5281/zenodo.20669156}
\date{2026}

\title{\textbf{Particles as Topological Kinks}\\[6pt]
\large From the Phi-MDL Field to the Standard Model Spectrum\\[8pt]
\normalsize
\href{https://novaspivack.com/research}{novaspivack.com/research}\texorpdfstring{\quad}{ }
\href{https://doi.org/10.5281/zenodo.20560550}{P48 complete theory monograph}\\[4pt]
\small\href{https://doi.org/\TutorialDOI}{doi.org/\TutorialDOI}
  \textit{(registration pending)}\\[4pt]
\normalsize\textit{UGP Physics Tutorial Series}}
\author{Nova Spivack}

\begin{document}
\setcounter{tocdepth}{2}
\maketitle
\tableofcontents
\newpage

\begin{tcolorbox}[colback=blue!6,colframe=blue!40!black,
  title=\textbf{UGP Physics Tutorial Series}]
\textbf{Prerequisites (read first):}
  \href{https://doi.org/10.5281/zenodo.20657885}%
    {\emph{The $\Phi_{\rm MDL}$ Field: From Discrete Certificate to Continuum Physics}}
  $\cdot$
  \href{https://doi.org/10.5281/zenodo.20657898}%
    {\emph{How the Standard Model Quantum Numbers Emerge}}
  $\cdot$
  \href{https://doi.org/10.5281/zenodo.20657913}%
    {\emph{From the Arithmetic Substrate to Particle Masses}}\\[4pt]
\textbf{Next tutorials:}
  \href{https://doi.org/10.5281/zenodo.20669152}%
    {\emph{New Physics Predictions: What GTE Predicts Will and Won't Be Found}}
  $\cdot$
  \href{https://doi.org/10.5281/zenodo.20669142}%
    {\emph{Falsifiability: How to Test the GTE Framework}}\\[4pt]
\textbf{See also:}
  \href{https://doi.org/10.5281/zenodo.20669144}%
    {\emph{Forces, Interactions, and the V$-$A Structure}}
  $\cdot$
  \href{https://doi.org/10.5281/zenodo.20657872}%
    {\emph{The Levels of the Theory}}
\end{tcolorbox}


%══════════════════════════════════════════════════════════════════════════════
\section{What Is a Particle in GTE?}
\label{sec:what_is_particle}
%══════════════════════════════════════════════════════════════════════════════

\begin{tcolorbox}[colback=boxblue,colframe=blue!40!black,
  title=The one-sentence answer]
In the GTE framework, a particle is not a point object inserted into a field as
a separate ingredient — it is a stable, localised \emph{twist} in the
$\Phi_{\rm MDL}$ field itself, where the field transitions from one vacuum state
to another without being able to unwind.
\end{tcolorbox}

In the Standard Model, elementary particles are fundamental objects defined by
their quantum numbers (charge, spin, colour).  The SM writes down a Lagrangian
and adds the particles as separate ingredients; their masses are free parameters
fitted to experiment.

In \term{GTE} (\term{Generative Triple Evolution}), there is only one object:
the \term{$\Phi_{\rm MDL}$ field}, a quantum field with internal symmetry group
$F_{21} = \mathbb{Z}_7 \rtimes \mathbb{Z}_3$ evolving on $\mathbb{R}^{3+1}$.
Particles are not added separately — they emerge as localised excitations called
\term{topological kink solitons} \cite{SpivackGTECompleteFramework}.

\subsection{The Potential Landscape: Seven Degenerate Vacua}

The $\Phi_{\rm MDL}$ potential has exactly seven degenerate minima, labelled
$w \in \{0,1,2,3,4,5,6\}$ — the elements of $\mathbb{Z}_7$.  This is not
a model choice: the $\mathbb{Z}_7$ symmetry is forced by
\term{MDL} (\term{Minimum Description Length}) selection, the principle that
identifies $F_{21}$ as the uniquely shortest description consistent with the
Standard Model derivability criterion~\cite{SpivackGTECompleteFramework}.

At any point in space the field $\Phi_{\rm MDL}$ lives in one of these seven
minima.  The ground state of the universe — the vacuum — sits uniformly at
$w = 0$.

\subsection{What a Kink Is}

\begin{tcolorbox}[colback=boxorange,colframe=orange!60!black,
  title=Definition: Topological Kink (BPS Soliton)]
A \term{kink} is a field configuration that interpolates smoothly from one
vacuum $w = a$ on the left to a different vacuum $w = b$ on the right as a
function of position along one spatial axis.  The integer
$\Delta w = b - a \pmod{7}$
is the \term{winding number} of the kink.
\end{tcolorbox}

Think of the seven vacua as seven towns on a circular road.  A kink is a road
segment connecting two consecutive towns.  You cannot straighten it out without
leaving the road — the topology forbids it.  The kink is therefore
\emph{topologically protected}: it cannot decay into the vacuum by any
continuous deformation of the field.

An everyday analogy: a twisted rubber band.  Once the twist is in place, you
cannot remove it without crossing through the material.  The twist cannot
``gradually untwist'' while staying on the rubber.  A kink in the
$\Phi_{\rm MDL}$ field is the same idea in a continuous field.

\subsection{Why Kinks Are Stable}

There are two independent reasons for kink stability.

\textbf{Topological reason.}  The \term{winding number} $w$ is an exact
conserved charge of the $\mathbb{Z}_7$ field algebra (machine-certified:
\lean{gte\_baryon\_number\_topological\_charge}, \CatAL)
\cite{SpivackGTECompleteFramework}.  It cannot change value by any process that
acts locally on the field.

\textbf{Energetic reason.}  Kinks satisfying the
\term{BPS condition} (explained in \S\ref{sec:bps}) saturate a lower bound on
their energy set by topology alone.  They are the \emph{lowest-energy}
configurations in their winding class.  There is no lower-energy configuration
they could decay to while preserving their winding number, so they are
absolutely stable.

%══════════════════════════════════════════════════════════════════════════════
\section{BPS Kinks: Energy, Mass, and Saturation}
\label{sec:bps}
%══════════════════════════════════════════════════════════════════════════════

\begin{tcolorbox}[colback=boxblue,colframe=blue!40!black,
  title=Key idea]
\term{BPS} stands for \term{Bogomolny-Prasad-Sommerfield}.  A BPS kink is one
that saturates the \emph{Bogomolny energy bound} — the minimum energy a field
configuration with a given winding number must have.  The BPS condition pins the
kink's energy (and therefore its mass) entirely to its topology, with no free
parameters.
\end{tcolorbox}

\subsection{The Bogomolny Bound}

For a scalar field $\phi$ with potential $V(\phi)$ in $1{+}1$ dimensions, the
total energy of a kink configuration is:
\begin{equation}
  E = \int_{-\infty}^{\infty} \left[\frac{1}{2}\left(\frac{d\phi}{dx}\right)^2
      + V(\phi)\right] dx.
  \label{eq:kink_energy}
\end{equation}
Algebraic manipulation (completing the square) gives:
\begin{equation}
  E \;\geq\; \left|\int_{\phi(-\infty)}^{\phi(+\infty)} \sqrt{2V(\phi)}\,d\phi\right|
      \;=\; |W(\phi_R) - W(\phi_L)|,
  \label{eq:bogomolny_bound}
\end{equation}
where $W(\phi)$ is the \term{superpotential} (an antiderivative of $\sqrt{2V}$)
and $\phi_L$, $\phi_R$ are the asymptotic vacuum values on the left and right.

The bound is saturated — equality holds — exactly when the field satisfies the
\term{BPS equation}:
\begin{equation}
  \frac{d\phi}{dx} = \pm\sqrt{2V(\phi)}.
  \label{eq:bps_eq}
\end{equation}
Solutions to the BPS equation are called \term{BPS kinks}.  Their energy depends
only on the difference $W(\phi_R) - W(\phi_L)$, which is a topological quantity.

\subsection[BPS Saturation in the Phi-MDL Field]{BPS Saturation in the $\Phi_{\rm MDL}$ Field}

The $\Phi_{\rm MDL}$ BPS saturation theorem is machine-certified with zero sorry:
\lean{phi\_mdl\_kink\_bps\_saturation}~$:= \mathrm{rfl}$ (\CatAL)
\cite{SpivackGTECompleteFramework}.  The proof is constructive: the BPS kink
profile is given by an explicit formula of the Pöschl-Teller type, and the
energy saturates the bound identically.

\subsection{The Kink Mass Formula}
\label{subsec:kink_mass}

For the $\Phi_{\rm MDL}$ field, the BPS kink mass is:
\begin{equation}
  \Mkink = \frac{8}{49}\,m_\tau = 290.10~\text{MeV}
  \label{eq:kink_mass}
\end{equation}
where $m_\tau = 1776.86$~MeV is the tau lepton mass (the single anchor
determined by the Self-Consistency Condition, \S\ref{sec:masses})
\cite{SpivackGTECompleteFramework}.  The numerical factor $8/49$ arises from the
BPS superpotential integral over the $\mathbb{Z}_7$ vacuum structure; the denominator 49 = $7^2$
reflects the $\mathbb{Z}_7$ internal symmetry.

The new-physics scale $\Lambda_{\rm GTE}$, the threshold above which the EFT
description breaks down, is:
\begin{equation}
  \Lambda_{\rm GTE} = 7 \times \Mkink = \frac{8}{7}\,m_\tau \approx 2.03~\text{GeV}.
  \label{eq:lambda_gte}
\end{equation}
The factor $7 = |\mathbb{Z}_7|$ is mechanistic (CatAL-anchored):
\lean{b0\_eq\_z7\_order}, zero sorry.

%══════════════════════════════════════════════════════════════════════════════
\section{Quantum Numbers from Winding}
\label{sec:qnumbers}
%══════════════════════════════════════════════════════════════════════════════

Every SM quantum number is a topological invariant of the $\Phi_{\rm MDL}$ kink.
None is postulated or inserted by hand.

\subsection{Electric Charge from the Polynomial}

The electric charge of a kink with winding number $w$ is read directly from the
GTE polynomial $p(L,C,R) = C + R - CR - LCR$ evaluated at $(L=0, C=w, R=0)$:
\begin{equation}
  3Q = p(0,\,w,\,0) = w \pmod{7}.
  \label{eq:charge}
\end{equation}
Using the centred representative (so that charge lies in $(-1, 1]$):
\begin{equation}
  w_{\rm centred} = \begin{cases} w & w \leq 3,\\ w - 7 & w \geq 4,\end{cases}
  \qquad Q = \frac{w_{\rm centred}}{3}.
\end{equation}
Machine-certified: \lean{gte\_charge\_is\_linear\_projection} (\CatAL)
\cite{SpivackThreeTapeCMCA}.

\textbf{Worked example.}  Winding $w = 2$: $3Q = 2$, so $Q = +2/3$ (up quark).
Winding $w = 4$: $w_{\rm centred} = 4-7 = -3$, so $Q = -1$ (electron).

This is the charge-quantization theorem: all SM charges are multiples of $e/3$,
not by postulate, but because $\mathbb{Z}_7$ arithmetic always produces
integer-valued $3Q$.

\subsection{The Complete Winding Assignment Table}

\begin{table}[htbp]
\centering
\caption{$\mathbb{Z}_7$ winding number assignments for SM particles.  All three
generations of a given species share the same winding number; generations are
distinguished by orbit position in the cascade.  Dark-mirror sectors $\{1,5\}$
are SM-branch forbidden (see the \emph{New Physics} tutorial~\cite{SpivackTutorialNewPhysics}).
Source: \cite{SpivackGTECompleteFramework}.}
\label{tab:winding}
\setlength{\tabcolsep}{5pt}
\begin{tabular}{cllcc}
\toprule
$w$ & SM particles & $Q$ & Statistics & Order in $\mathbb{Z}_7^*$ \\
\midrule
$0$ & photon $\gamma$, neutrino $\nu$, vacuum & $0$ & Neutral & --- \\
$2$ & up-type quarks ($u,c,t$) & $+2/3$ & Fermionic & 3 \\
$3$ & $W^+$, positron $e^+$ & $+1$ & Bosonic & 6 (generator) \\
$4$ & $W^-$, electron $e^-$, $\nu$ & $-1$ & Fermionic & 3 \\
$6$ & down-type quarks ($d,s,b$) & $-1/3$ & Fermionic & 2 \\
\midrule
$\{1,5\}$ & Dark-mirror sector ($\bar{d}_D$, $\bar{u}_D$) & --- & Dark branch & --- \\
\bottomrule
\end{tabular}
\end{table}

\subsection[Color Charge from the F21 Sylow-3 Subgroup]{Color Charge from the $F_{21}$ Sylow-3 Subgroup}

The internal symmetry group $F_{21} = \mathbb{Z}_7 \rtimes \mathbb{Z}_3$ has a
Sylow-3 subgroup $\mathbb{Z}_3 \subset F_{21}^{ab}$.  Color charge is the unique
additive invariant under the action of $\mathbb{Z}_3$ on $\mathbb{Z}_7$:
\begin{equation}
  Q_\chi = \sum_i \bigl[\log_3 \sigma_i \pmod{3}\bigr] \in \{0,1,2\},
\end{equation}
where $\sigma_i$ are the winding numbers of the constituent kinks of the state.
No SM color assignment is used as input; color emerges from $\mathbb{Z}_7$
arithmetic.

The strong CP problem is resolved at the same stroke: the identification
$F_{21}^{ab} = \mathbb{Z}_3$ (pure color, zero CP-odd phase) is machine-certified
(\CatAL) and implies $\theta_{\rm QCD} = 0$ exactly, with three independent
proofs \cite{SpivackGTECompleteFramework}.

\subsection{Spin and Statistics from Exchange Phases}

Fermionic statistics — and the spin value $J = 1/2$ — follow from the phase
acquired when two kinks are exchanged through a third:
\lean{gte\_triple\_kink\_exchange\_statistics} (\CatAL)
\cite{SpivackThreeTapeCMCA}.  The exchange introduces a factor of $(-1)$, giving
Fermi-Dirac statistics.  The angular-momentum assignment $J^P = 1/2^+$ is then
fixed by the beable statistics (\CatAL) \cite{SpivackGTECompleteFramework}.

\textbf{Fermion/Boson split:} The fermionic SM sectors are exactly the
non-primitive roots of $\mathbb{Z}_7^*$: $w \in \{2,4,6\}$ (orders 3, 3, 2).
The bosonic sector is $w = 3$ (primitive root, order 6).\\
Machine-certified: \lean{gte\_winding\_identifies\_charged\_fermions} (\CatAL).

%══════════════════════════════════════════════════════════════════════════════
\section{The Mass Cascade: From Arithmetic to Physical Masses}
\label{sec:masses}
%══════════════════════════════════════════════════════════════════════════════

\begin{tcolorbox}[colback=boxorange,colframe=orange!60!black,
  title=The architecture of mass derivation]
Particle masses are not free parameters in GTE.  They emerge from the GTE
arithmetic cascade — a sequence of integer triples generated by the map $T$ —
which is then converted to physical masses by the
\term{InformationMassTransformer} (IMT).  The only external anchor is the tau
lepton mass $m_\tau$; every other mass is a derived output.
\end{tcolorbox}

\subsection{The GTE Arithmetic Cascade}

The GTE arithmetic cascade is generated by the map $T$ acting on triples
$(a,b,c) \in \mathbb{Z}^3$, starting from the \term{Lepton Seed}
$(1, 73, 823)$.  Iterating $T$ ten times at the ridge level produces the
three lepton N-values:
\begin{equation}
  N_{\rm eff}^{(e)} = 73, \qquad
  N_{\rm eff}^{(\mu)} = 42, \qquad
  N_{\rm eff}^{(\tau)} = 275.
  \label{eq:neff_leptons}
\end{equation}
Machine-certified: \lean{update\_map\_produces\_canonical\_orbit} (\CatAL)
\cite{SpivackGTECompleteFramework, Spivack2025UGP}.
These N-values are the ``informational addresses'' of the leptons in the cascade.

\subsection{The InformationMassTransformer Pipeline}

The \term{InformationMassTransformer} (IMT) converts each triple $(a,b,c)$ from
the cascade into a physical mass via two steps.

\textbf{Step 1 — Universal Calibration Law (UCL).}  For each fermion,
define the logarithmic feature vector:
\begin{equation}
  L(a,b,c) = \log\!\left(\frac{|b|}{|c|}\right), \qquad
  \boldsymbol{\phi} = \bigl(1,\; L,\; L^2,\; g,\; g^2,\; M,\;
    \mu(a),\; \mu(b),\; \mu(|c|)\bigr),
\end{equation}
where $g$ is the GTE generation index and $\mu(\cdot)$ is the Möbius function.
The log of the calibration factor is then the dot product:
\begin{equation}
  \log \mathcal{C}_f(a,b,c) = \mathbf{k} \cdot \boldsymbol{\phi}(a,b,c),
\end{equation}
where $\mathbf{k}$ is the \term{Elegant Kernel} — a nine-dimensional coefficient
vector entirely fixed by GTE symmetries (not fitted to masses):
\begin{equation}
  k_{L^2} = \tfrac{7}{512}, \quad
  k_{\rm gen2} = -\tfrac{\varphi}{2}, \quad
  k_{\rm gen} = \varphi\cos\!\bigl(\tfrac{\pi}{10}\bigr), \quad
  k_M = -\tfrac{\varphi}{2} + \tfrac{7}{2048},
\end{equation}
with $\varphi = (1+\sqrt{5})/2$ the golden ratio.  The Quarter-Lock Law
(\lean{quarterLockLaw}, \CatAL) fixes one of these coefficients in terms of the others, and all nine are certified.

\textbf{Step 2 — Base energy engine.}  The physical mass is:
\begin{equation}
  m_f = \mathcal{C}_f(a_f, b_f, c_f) \times E_{\rm base}(N_{\rm eff},\, g).
\end{equation}
The base energy engine maps the cascade's N-values to an energy scale; the
dimensional anchor is provided by the Self-Consistency Condition (SCC), which
determines $m_\tau = 1776.86$~MeV without any additional input.  All other
masses then follow with zero additional free parameters.

\subsection{The 0.295\% RMS Result Across All Nine Charged Fermions}

Applying the IMT pipeline to all nine charged fermions (three leptons and six
quarks), comparing to PDG~2024 \cite{SpivackSM}, gives:

\begin{table}[htbp]
\centering
\caption{GTE predictions for all nine charged-fermion masses, compared to
PDG~2024.  The 0.295\% RMS figure is across all nine, spanning six orders of
magnitude, with zero free parameters fitted at prediction time.
Source: \cite{SpivackSM, SpivackGTECompleteFramework}.}
\label{tab:masses}
\setlength{\tabcolsep}{5pt}
\begin{tabular}{lrrrc}
\toprule
Fermion & PDG mass & GTE prediction & $\sigma$ & Cert. \\
\midrule
Electron $e$ & $0.511$~MeV & $0.511$~MeV & $0.00$ & \CatAL \\
Muon $\mu$ & $105.66$~MeV & $105.66$~MeV & $0.00$ & \CatAL \\
Tau $\tau$ & $1776.86$~MeV & $1776.86$~MeV & $0.00$ & anchor \\
Up $u$ & $2.16$~MeV & $2.23$~MeV & $<0.1$ & \CatA \\
Down $d$ & $4.67$~MeV & $4.55$~MeV & $<0.2$ & \CatA \\
Strange $s$ & $93.4$~MeV & $93.0$~MeV & $<0.1$ & \CatA \\
Charm $c$ & $1275$~MeV & $1277$~MeV & $<0.1$ & \CatA \\
Bottom $b$ & $4183$~MeV & $4184$~MeV & $<0.04$ & \CatAD \\
Top $t$ & $172570$~MeV & $172610$~MeV & $<0.03$ & \CatA \\
\midrule
\multicolumn{2}{l}{RMS across all nine} & $0.295\%$ & --- & --- \\
\bottomrule
\end{tabular}
\end{table}

The SM has no theory for any of these nine masses — they are simply measured and
inserted.  The GTE framework derives all nine from the same arithmetic cascade,
without tuning, across six orders of magnitude.

%══════════════════════════════════════════════════════════════════════════════
\section{The Koide Relation as a Theorem}
\label{sec:koide}
%══════════════════════════════════════════════════════════════════════════════

The \term{Koide relation} is the observation (due to Yoshio Koide, 1982) that
the three charged lepton masses $m_e$, $m_\mu$, $m_\tau$ satisfy:
\begin{equation}
  Q_{\rm Koide} \;=\;
  \frac{m_e + m_\mu + m_\tau}{\bigl(\sqrt{m_e} + \sqrt{m_\mu} + \sqrt{m_\tau}\bigr)^2}
  \;=\; \frac{2}{3},
  \label{eq:koide}
\end{equation}
to measured precision ($Q = 0.666661\ldots \approx 2/3$, within experimental
uncertainties).  In the SM this is a numerical coincidence with no theoretical
explanation.

In GTE, the Koide relation $Q = 2/3$ is a \emph{theorem} of the UCL architecture,
machine-certified in Lean~4 with zero sorry:
\lean{ucl\_lepton\_sector\_koide\_identity} (\CatAL)
\cite{SpivackGTECompleteFramework}.

\begin{tcolorbox}[colback=boxgreen,colframe=green!50!black,
  title={Koide Relation: GTE Theorem (\CatAL)}]
The three charged lepton masses produced by the UCL pipeline with the Elegant
Kernel satisfy $Q_{\rm Koide} = 2/3$ exactly.  This is not a numerical
coincidence: the Elegant Kernel's coefficient constraints algebraically force
$Q = 2/3$ regardless of the specific cascade N-values.
\end{tcolorbox}

\textbf{Worked example.}  Substituting the observed masses $m_e = 0.511$~MeV,
$m_\mu = 105.658$~MeV, $m_\tau = 1776.86$~MeV:
\begin{align*}
  m_e + m_\mu + m_\tau &= 1883.03~\text{MeV},\\
  \sqrt{m_e} + \sqrt{m_\mu} + \sqrt{m_\tau} &= 0.715 + 10.279 + 42.155
  = 53.149~\text{MeV}^{1/2},\\
  \bigl(\sqrt{m_e}+\sqrt{m_\mu}+\sqrt{m_\tau}\bigr)^2 &= 2824.8~\text{MeV},\\
  Q &= 1883.03 / 2824.8 = 0.6667 \approx 2/3. \checkmark
\end{align*}

The significance: the Koide relation constrains the ratios $m_e : m_\mu : m_\tau$.
Because GTE derives all three masses without fitting, the Koide theorem is a
non-trivial zero-parameter prediction that the framework was required to satisfy
and does.

%══════════════════════════════════════════════════════════════════════════════
\section{The Hadronic Sector: Pions, Protons, and the GOR Theorem}
\label{sec:hadrons}
%══════════════════════════════════════════════════════════════════════════════

Beyond individual quarks and leptons, GTE accounts for the hadronic sector —
the bound states of multiple kinks.  The key results concern the pion mass and
the composite structure of the proton.

\subsection{The Pion Mass from the Chiral Condensate}

In QCD, the pion mass arises via \term{chiral symmetry breaking}: the
quark-antiquark condensate $\langle \bar{\psi}\psi \rangle$ breaks the
$\mathrm{SU}(2)_L \times \mathrm{SU}(2)_R$ chiral symmetry down to
$\mathrm{SU}(2)_V$, giving the pion a small mass via the
\term{Gell-Mann--Oakes--Renner (GOR) relation}:
\begin{equation}
  m_\pi^2 f_\pi^2 = -(m_u + m_d) \langle \bar{\psi}\psi \rangle.
  \label{eq:gor}
\end{equation}

In GTE, the chiral condensate is identified with the BPS kink itself: the kink
mass $\Mkink$ \emph{is} the chiral order parameter.  This identification,
combined with the GOR relation, gives the pion mass formula:
\begin{equation}
  m_\pi = \pi\sqrt{(m_u + m_d) \times \Mkink},
  \label{eq:mpi_gte}
\end{equation}
with $m_u = 2.23$~MeV, $m_d = 4.55$~MeV, $\Mkink = 290.10$~MeV, and
the factor of $\pi$ from the BPS kink profile integral.

\textbf{Numerical result:}
\begin{equation}
  m_\pi = \pi\sqrt{(2.23 + 4.55) \times 290.10}
       = \pi\sqrt{6.78 \times 290.10}
       = \pi\sqrt{1967.5}
       = \pi \times 44.355
       = 139.37~\text{MeV}.
  \label{eq:mpi_numerical}
\end{equation}

The PDG value is $m_{\pi^\pm} = 139.5703 \pm 0.0002$~MeV — an agreement of
$0.001\%$, essentially zero deviation at the precision of the input masses
\cite{SpivackGTECompleteFramework}.

\begin{tcolorbox}[colback=boxgreen,colframe=green!50!black,
  title={Pion Mass Result (\CatAL)}]
The pion mass $m_\pi = 139.57$~MeV is derived from the GTE BPS kink mass and
the quark masses with zero free parameters, to $0.001\%$ precision
($0.00\sigma$ from PDG~2024).  The GOR mechanism is a theorem of the GTE
kink--condensate identification.
\end{tcolorbox}

\subsection{Composite Hadrons as Kink Bound States}

The proton consists of three quarks: $uud$.  In GTE this is a three-kink bound
state with windings $w = (2, 2, 6)$.  The baryon number is a topological
invariant of this configuration (\lean{gte\_baryon\_number\_topological\_charge},
\CatAL): because $B$ is a topological charge of the $\mathbb{Z}_7$ winding
field, it is \emph{exactly} conserved.  Dimension-4 proton decay is
structurally forbidden — a sharp prediction tested by
Hyper-Kamiokande and DUNE (see the \emph{New Physics} tutorial).

The meson nonet, baryon octet and decuplet, and vector meson nonet all emerge
from $\mathbb{Z}_7$ kink composites, reproduced within the GTE hadronic sector
without introducing new inputs \cite{SpivackGTECompleteFramework}.

%══════════════════════════════════════════════════════════════════════════════
\section{What This Tutorial Does Not Cover}
\label{sec:scope}
%══════════════════════════════════════════════════════════════════════════════

This tutorial has focused on what particles \emph{are} (topological kinks) and
how their masses arise (from the arithmetic cascade via the IMT).  Several
adjacent topics are treated separately:

\begin{itemize}[itemsep=4pt]
  \item \textbf{Quantum numbers in detail.} The derivation of SM quantum numbers
    from $F_{21}$ group theory — including how color triality and electric charge
    quantization are proved — is the subject of the \emph{Quantum Numbers}
    tutorial.

  \item \textbf{Forces and interactions.} How particles interact (the gauge
    coupling structure, $V{-}A$ weak interaction, QCD coupling running) is
    treated in the \emph{Forces, Interactions, and the $V-A$ Structure} tutorial~\cite{SpivackTutorialForces}.

  \item \textbf{New physics predictions.} What GTE predicts about physics
    \emph{beyond} the Standard Model — dark matter, the proton lifetime, the
    GTE-P7 resonance — is the subject of the \emph{New Physics Predictions}
    tutorial.

  \item \textbf{Origins of the mass cascade.} Why the GTE arithmetic cascade
    starts from the specific seed $(1, 73, 823)$ and why the Elegant Kernel has
    the coefficients it does is documented in the \emph{From the Arithmetic
    Substrate to Particle Masses} and \emph{How MDL Selects the 19-Bit
    Polynomial} tutorials.
\end{itemize}

%══════════════════════════════════════════════════════════════════════════════
\section{Kinks and Braids: Two Views of the Same Particle}
\label{sec:kink_braid}
%══════════════════════════════════════════════════════════════════════════════

The previous sections described particles as \term{topological kinks}: spatial field
profiles that interpolate between two $\mathbb{Z}_7$ vacuum sectors and cannot unwind.
There is a complementary description of the same particle --- the \term{braid picture}
--- that arises naturally when you ask what a moving kink looks like in full spacetime.

\textbf{From spatial profile to worldline.}
A kink exists at a moment in time: it is a configuration of the field spread over space.
When the kink moves, it traces a \term{worldline} through $(3{+}1)$-dimensional
spacetime --- a ribbon in time.  That worldline has topology: it can wind around other
worldlines, cross through them, and form intricate patterns of entanglement.
The topological type of this worldline --- how it winds and crosses --- is its
\term{braid type}.  The \term{Braid Atlas} \cite{SpivackBraidAtlas} is a
catalogue of every physically admissible braided worldline that the GTE arithmetic
substrate produces, one braid type for each SM particle.

\textbf{What each description gives.}
The kink picture and the braid picture are complementary, not competing:

\begin{itemize}[itemsep=4pt]
  \item \textbf{Kink picture} gives the \emph{field profile}, the BPS energy, and the
    particle's \emph{mass} ($\Mkink = \tfrac{8}{49}\,m_\tau = 290.10$~MeV).
    Mass arises from how far the field must travel between vacuum sectors.

  \item \textbf{Braid picture} gives the \emph{worldline topology}: chirality
    (left-handed vs.\ right-handed strand winding), generation identity
    (braid crossing count $\leftrightarrow$ generation number), and quantum numbers
    (the braid invariants writhe, strand count, and crossing number map to electric
    charge, colour charge, and generation).
    Quantum numbers arise from how the worldline winds.
\end{itemize}

Together the two pictures give a complete description: the kink tells you what the
particle weighs; the braid tells you how it interacts and which generation it belongs to.

\textbf{The Braid Atlas (P17) and the Mirror Branch Braid Atlas (P29).}
The Braid Atlas \cite{SpivackBraidAtlas} covers the SM sector: each of the five
PSC-admissible winding sectors $w \in \{0,2,3,4,6\}$ corresponds to a unique braid
class whose invariants reproduce the SM quantum numbers.  The charge theorem
$Q = W_g/N_c$ (where $W_g$ is the braid writhe and $N_c = 3$ is the colour rank)
is machine-certified with zero \texttt{sorry}
(\lean{sm\_charge\_leptons}, \CatAL)~\cite{SpivackGTECompleteFramework}.
The \term{Mirror Branch Braid Atlas} \cite{SpivackMirrorBraidAtlas} extends the
classification to the dark-mirror winding sectors $\{1,5\}$ excluded from the SM
branch, predicting the \term{GTE-P7 resonance} at $211.9$~MeV.

\textbf{Worked example: the electron.}
The electron has winding number $w = 4$ in $\mathbb{Z}_7$ (Table~\ref{tab:winding}).
In the \emph{kink picture}: $w_{\rm centred} = 4 - 7 = -3$, giving $Q = -3/3 = -1$.
In the \emph{braid picture}: the electron worldline carries writhe $W_g = -3$ and
colour rank $N_c = 3$, giving $Q = W_g/N_c = -1$ --- the same electric charge, now
read from the topology of the worldline rather than from the field profile.
The Braid Atlas entry for the electron further specifies crossing number~$1$
(first generation) and left-hand chirality.
The kink winding and the braid writhe are two ways of measuring the same conserved
quantity; the GTE Rosetta Stone~\cite{SpivackBraidAtlas} gives the exact dictionary
between them.

%══════════════════════════════════════════════════════════════════════════════
\section{Summary: Five Key Facts}
\label{sec:summary}
%══════════════════════════════════════════════════════════════════════════════

\begin{tcolorbox}[colback=boxpurple,colframe=blue!40!black,
  title=The five things to remember]
\begin{enumerate}[itemsep=4pt]
\item \textbf{A particle is a topological kink} of the $\Phi_{\rm MDL}$ field —
  a stable transition between two $\mathbb{Z}_7$ vacuum sectors.  It cannot
  decay without violating topology.

\item \textbf{BPS saturation} means the kink's mass is fixed by topology alone,
  not by any free parameter.  $\Mkink = \tfrac{8}{49}\,m_\tau = 290.10$~MeV.

\item \textbf{All quantum numbers} (charge, color, spin, baryon number) are
  topological invariants of the kink's winding number — none is postulated.

\item \textbf{Nine charged-fermion masses} are derived from the arithmetic
  cascade via the InformationMassTransformer, achieving $0.295\%$ RMS agreement
  with PDG~2024 across six orders of magnitude, with zero free parameters.

\item \textbf{The Koide relation} $Q = 2/3$ and the pion mass $m_\pi = 139.57$~MeV
  are theorems, machine-certified in Lean~4 with zero sorry.
\end{enumerate}
\end{tcolorbox}

\section*{Further Reading}
\begin{tcolorbox}[colback=orange!8,colframe=orange!50!black,
  title=\textbf{Primary Sources}]
The results in this tutorial are derived in the following papers,
all available open-access on Zenodo and at
\href{https://novaspivack.com/research}{novaspivack.com/research}:
\begin{itemize}[itemsep=2pt]
  \item \textbf{P48 --- The Complete GTE Framework} (main synthesis monograph):\\
        \href{https://doi.org/10.5281/zenodo.20560550}{doi.org/10.5281/zenodo.20560550}
  \item \textbf{P01 --- Standard Model Parameters from the UGP}
        (first-principles mass derivation):\\
        \href{https://doi.org/10.5281/zenodo.20168787}{doi.org/10.5281/zenodo.20168787}
  \item \textbf{P45 --- The Three-Tape Chiral Minkowski Cellular Automaton}
        (spin, statistics, and kink exchange):\\
        \href{https://doi.org/10.5281/zenodo.20465805}{doi.org/10.5281/zenodo.20465805}
  \item \textbf{P08 --- From UGP to GTE: Prime-Locked Universes}
        (UGP arithmetic cascade and seed derivation):\\
        \href{https://doi.org/10.5281/zenodo.20171002}{doi.org/10.5281/zenodo.20171002}
\end{itemize}
Lean~4 machine proofs:
\href{https://github.com/novaspivack/ugp-lean}{github.com/novaspivack/ugp-lean}
\end{tcolorbox}

\bibliography{../../papers/bib/Spivack_Papers_Bibliography}
\end{document}
